\SourceLeaf{056}
functionem solius $x$, eam, in quam transit $r$ per substitutiones $l'=L',l''=L''$,
$l'''=L'''$ etc., patet determinantem functionis $z$ fieri
\[
=p^{m-2}y^{(m-1)(m-2)}r
\]
determinantem functionis $Z$ autem
\[
=P^{m-2}Y^{(m-1)(m-2)}R
\]
Quare quum per hypothesin $P$ non sit $=0$, res iam in eo vertitur, ut demonstremus, $R$ certo identice evanescere non posse.

\SecondArticle{15}
Ad hunc finem adhuc aliam indeterminatam $w$ introducemus, atque productum ex omnibus
\[
(a+b-c-d)w+(a-c)(a-d)
\]
exclusis repetitionibus considerabimus, quod quum ipsas $a,b,c$ etc. symmetrice
involvat, tamquam functio integra indeterminatarum $w,\lambda',\lambda'',\lambda'''$ etc. exhiberi
poterit. Denotabimus hanc functionem per $f(w,\lambda',\lambda'',\lambda'''\text{ etc.})$. Multitudo illorum factorum $(a+b-c-d)w+(a-c)(a-d)$ erit
\[
=\frac12m(m-1)(m-2)(m-3)
\]
unde facile colligimus, fieri
\[
f(0,\lambda',\lambda'',\lambda'''\text{ etc.})=\pi^{(m-2)(m-3)}
\]
et proin etiam
\[
f(0,l',l'',l'''\text{ etc.})=p^{(m-2)(m-3)}
\]
nec non
\[
f(0,L',L'',L'''\text{ etc.})=P^{(m-2)(m-3)}
\]
Functio $f(w,\allowbreak L',\allowbreak L'',\allowbreak L'''\text{ etc.})$ generaliter quidem loquendo ad ordinem
\[
\frac12m(m-1)(m-2)(m-3)
\]
referenda erit: at in casibus specialibus utique ad ordinem inferiorem pertinere
potest, si forte contingat, ut quidam coëfficientes inde ab altissima potestate ipsius
$w$ evanescant: impossibile autem est, ut illa functio tota sit identice $=0$, quum
\SourceLeaf{057}
aequatio modo inventa doceat, functionis saltem terminum ultimum non evanescere. Supponemus, terminum altissimum functionis $f(w,\allowbreak L',\allowbreak L'',\allowbreak L'''\text{ etc.})$, qui
quidem coëfficientem non evanescentem habeat, esse $Nw^\nu$. Si igitur substituimus
$w=x-a$, patet, $f(x-a,L',L'',L'''\text{ etc.})$ esse functionem integram indeterminatarum $x,a$, sive quod idem est, functionem ipsius $x$ cum coëfficientibus ab
indeterminata $a$ pendentibus, ita tamen ut terminus altissimus sit $Nx^\nu$, et proin
coëfficientem determinatum ab $a$ non pendentem habeat, qui non sit $=0$. Perinde $f(x-b,L',L'',L'''\text{ etc.}),f(x-c,L',L'',L'''\text{ etc.})$ erunt functiones integrae
indeterminatae $x$, tales ut singularum terminus altissimus sit $Nx^\nu$, terminorum
sequentium autem coëfficientes resp. a $b,c$ etc. pendeant. Hinc productum ex
$m$ factoribus
\[
\begin{gathered}
f(x-a,L',L'',L'''\text{ etc.})\\
f(x-b,L',L'',L'''\text{ etc.})\\
f(x-c,L',L'',L'''\text{ etc.})\\
\text{etc.}
\end{gathered}
\]
erit functio integra ipsius $x$, cuius terminus altissimus erit $N^mx^{m\nu}$, dum terminorum sequentium coëfficientes pendent ab indeterminatis $a,b,c$ etc.

Consideremus iam porro productum ex $m$ factoribus his
\[
\begin{gathered}
f(x-a,l',l'',l'''\text{ etc.})\\
f(x-b,l',l'',l'''\text{ etc.})\\
f(x-c,l',l'',l'''\text{ etc.})\\
\text{etc.}
\end{gathered}
\]
quod quum sit functio indeterminatarum, $x,a,b,c$ etc., $l',l'',l'''$ etc. et quidem
symmetrica respectu ipsarum $a,b,c$ etc., exhiberi poterit tamquam functio indeterminatarum $x,\lambda',\lambda'',\lambda'''$ etc. $l',l'',l'''$ etc. per
\[
\varphi(x,\lambda',\lambda'',\lambda'''\text{ etc.},\ l',l'',l'''\text{ etc.})
\]
denotanda. Erit itaque
\[
\varphi(x,\lambda',\lambda'',\lambda'''\text{ etc.},\ \lambda',\lambda'',\lambda'''\text{ etc.})
\]
productum ex factoribus
\[
\begin{gathered}
f(x-a,\lambda',\lambda'',\lambda'''\text{ etc.})\\
f(x-b,\lambda',\lambda'',\lambda'''\text{ etc.})\\
f(x-c,\lambda',\lambda'',\lambda'''\text{ etc.})\\
\text{etc.}
\end{gathered}
\]
\SourceLeaf{058}
et proin indefinite divisibilis per $\rho$, quum facile perspiciatur, quemlibet factorem
ipsius $\rho$ in aliquo illorum factorum implicari. Statuemus itaque
\[
\begin{aligned}
&\varphi(x,\lambda',\lambda'',\lambda'''\text{ etc.},\ \lambda',\lambda'',\lambda'''\text{ etc.})\\
&=\rho\psi(x,\lambda',\lambda'',\lambda'''\text{ etc.})
\end{aligned}
\]
ubi characteristica $\psi$ functionem integram exhibebit. Hinc vero facile deducitur,
etiam identice esse
\[
\begin{aligned}
&\varphi(x,\allowbreak L',\allowbreak L'',\allowbreak L'''\text{ etc.},\allowbreak \ L',\allowbreak L'',\allowbreak L'''\text{ etc.})\\
&=R\psi(x,L',L'',L'''\text{ etc.})
\end{aligned}
\]
Sed supra demonstravimus, productum e factoribus
\[
\begin{gathered}
f(x-a,L',L'',L'''\text{ etc.})\\
f(x-b,L',L'',L'''\text{ etc.})\\
f(x-c,L',L'',L'''\text{ etc.})\\
\text{etc.}
\end{gathered}
\]
quod erit $=\varphi(x,\lambda',\lambda'',\lambda'''\text{ etc.},\ L',L'',L'''\text{ etc.})$ habere terminum altissimum
$N^mx^{m\nu}$; eundem proin terminum altissimum habebit functio $\varphi(x,\allowbreak L',\allowbreak L'',\allowbreak L'''\text{ etc.},\allowbreak $
$L',L'',L'''\text{ etc.})$ adeoque certo non est identice $=0$. Quocirca etiam $R$ nequit
esse identice $=0$, neque adeo etiam determinans functionis $Z$. Q. E. D.

\SecondArticle{16}
\textsc{Theorema.} \textit{Denotet $\varphi(u,x)$}\HistoricalNote{12}\textit{ productum ex quotcunque factoribus talibus, in
quos indeterminatae $u,x$ lineariter tantum ingrediuntur, sive qui sint formae}
\[
\begin{gathered}
\alpha+\beta u+\gamma x\\
\alpha'+\beta'u+\gamma'x\\
\alpha''+\beta''u+\gamma''x\\
\text{etc.}
\end{gathered}
\]
\textit{sit porro $w$ alia indeterminata. Tunc functio}
\[
\varphi(u+w\, .\,\tfrac{\mathrm d\varphi(u,x)}{\mathrm dx},\ x-w\, .\,\tfrac{\mathrm d\varphi(u,x)}{\mathrm du})=\Omega
\]
\textit{indefinite erit divisibilis per $\varphi(u,x)$.}



\SourceLeaf{059}
\textit{Dem.} Statuendo
\[
\begin{aligned}
\varphi(u,x)&=(\alpha+\beta u+\gamma x)Q\\
&=(\alpha'+\beta'u+\gamma'x)Q'\\
&=(\alpha''+\beta''u+\gamma''x)Q''\\
&\quad\text{etc.}
\end{aligned}
\]
erunt $Q,Q',Q''$ etc. functiones integrae indeterminatarum $u,\allowbreak x,\allowbreak \alpha,\allowbreak \beta,\allowbreak \gamma,\allowbreak \alpha',\allowbreak \beta',\allowbreak \gamma'$,
$\alpha'',\beta'',\gamma''$ etc. atque
\[
\begin{aligned}
\frac{\mathrm d\varphi(u,x)}{\mathrm dx}&=\gamma Q+(\alpha+\beta u+\gamma x)\, .\,\frac{\mathrm dQ}{\mathrm dx}\\
&=\gamma'Q'+(\alpha'+\beta'u+\gamma'x)\, .\,\frac{\mathrm dQ'}{\mathrm dx}\\
&=\gamma''Q''+(\alpha''+\beta''u+\gamma''x)\, .\,\frac{\mathrm dQ''}{\mathrm dx}\\
&\quad\text{etc.}
\end{aligned}
\]
\[
\begin{aligned}
\frac{\mathrm d\varphi(u,x)}{\mathrm du}&=\beta Q+(\alpha+\beta u+\gamma x)\, .\,\frac{\mathrm dQ}{\mathrm du}\\
&=\beta'Q'+(\alpha'+\beta'u+\gamma'x)\, .\,\frac{\mathrm dQ'}{\mathrm du}\\
&=\beta''Q''+(\alpha''+\beta''u+\gamma''x)\, .\,\frac{\mathrm dQ''}{\mathrm du}\\
&\quad\text{etc.}
\end{aligned}
\]
Substitutis hisce valoribus in factoribus, e quibus conflatur productum $\Omega$,
puta in
\[
\begin{gathered}
\alpha+\beta u+\gamma x+\beta w\, .\,\frac{\mathrm d\varphi(u,x)}{\mathrm dx}-\gamma w\, .\,\frac{\mathrm d\varphi(u,x)}{\mathrm du}\\
\alpha'+\beta'u+\gamma'x+\beta'w\, .\,\frac{\mathrm d\varphi(u,x)}{\mathrm dx}-\gamma'w\, .\,\frac{\mathrm d\varphi(u,x)}{\mathrm du}\\
\alpha''+\beta''u+\gamma''x+\beta''w\, .\,\frac{\mathrm d\varphi(u,x)}{\mathrm dx}-\gamma''w\, .\,\frac{\mathrm d\varphi(u,x)}{\mathrm du}\\
\text{etc. resp.}
\end{gathered}
\]
hi obtinent valores sequentes
\[
\begin{gathered}
(\alpha+\beta u+\gamma x)(1+\beta w\, .\,\tfrac{\mathrm dQ}{\mathrm dx}-\gamma w\, .\,\tfrac{\mathrm dQ}{\mathrm du})\\
(\alpha'+\beta'u+\gamma'x)(1+\beta'w\, .\,\tfrac{\mathrm dQ'}{\mathrm dx}-\gamma'w\, .\,\tfrac{\mathrm dQ'}{\mathrm du})\\
(\alpha''+\beta''u+\gamma''x)(1+\beta''w\, .\,\tfrac{\mathrm dQ''}{\mathrm dx}-\gamma''w\, .\,\tfrac{\mathrm dQ''}{\mathrm du})\\
\text{etc.}
\end{gathered}
\]
